\textbf{Phần 1. Trắc nghiệm lựa chọn}
\setcounterpageref{ex}{0}

\begin{ex} %[5D7N37]
Trên bản đồ tỉ lệ $1 : 1\,000$, khoảng cách $1\text{ cm}$ trên bản đồ ứng với khoảng cách thật là
\choice
{$100\text{ cm}$}
{\True $1\,000\text{ cm}$}
{$10\text{ m}$}
{$100\text{ m}$}
\loigiai{
Khoảng cách thực tế là $1 \times 1\,000 = 1\,000\text{ cm}$.\\
\textbf{Đáp án của câu hỏi là B}
}
\end{ex}

\begin{ex} %[5D7N37]
Tỉ lệ bản đồ $1 : 500\,000$ cho biết
\choice
{\True Khoảng cách thực tế gấp $500\,000$ lần khoảng cách trên bản đồ}
{Khoảng cách trên bản đồ gấp $500\,000$ lần khoảng cách thực tế}
{Bản đồ thu nhỏ $50\,000$ lần}
{Độ dài thực tế bằng $500\,000\text{ m}$}
\loigiai{
Tỉ lệ $1 : 500\,000$ nghĩa là khoảng cách ngoài thực tế gấp $500\,000$ lần khoảng cách tương ứng vẽ trên bản đồ.\\
\textbf{Đáp án của câu hỏi là A}
}
\end{ex}

\begin{ex} %[5D7H37]
Trên bản đồ tỉ lệ $1 : 2\,000$, chiều dài một sân bóng đo được là $5\text{ cm}$. Chiều dài thật của sân bóng là
\choice
{$10\text{ m}$}
{$1\,000\text{ m}$}
{$50\text{ m}$}
{\True $100\text{ m}$}
\loigiai{
Chiều dài thật là: $5 \times 2\,000 = 10\,000\text{ cm} = 100\text{ m}$.\\
\textbf{Đáp án của câu hỏi là D}
}
\end{ex}

\begin{ex} %[5D7H37]
Khoảng cách giữa hai thành phố là $120\text{ km}$. Trên bản đồ tỉ lệ $1 : 2\,000\,000$, khoảng cách giữa hai thành phố đó là
\choice
{\True $6\text{ cm}$}
{$60\text{ cm}$}
{$12\text{ cm}$}
{$24\text{ cm}$}
\loigiai{
Đổi $120\text{ km} = 12\,000\,000\text{ cm}$.\\
Khoảng cách trên bản đồ là: $12\,000\,000 : 2\,000\,000 = 6\text{ (cm)}$.\\
\textbf{Đáp án của câu hỏi là A}
}
\end{ex}

\begin{ex} %[5D7H37]
Một thửa ruộng hình chữ nhật có chiều rộng thực tế $40\text{ m}$. Trên bản đồ tỉ lệ $1 : 1\,000$, chiều rộng của thửa ruộng đó là
\choice
{$40\text{ cm}$}
{$0{,}4\text{ cm}$}
{\True $4\text{ cm}$}
{$4\text{ mm}$}
\loigiai{
Đổi $40\text{ m} = 4\,000\text{ cm}$.\\
Khoảng cách trên bản đồ là: $4\,000 : 1\,000 = 4\text{ (cm)}$.\\
\textbf{Đáp án của câu hỏi là C}
}
\end{ex}

\begin{ex} %[5D7H37]
Độ dài thật của quãng đường từ nhà bạn An đến trường là $1{,}5\text{ km}$. Trên bản đồ, đoạn đường đó dài $3\text{ cm}$. Tỉ lệ của bản đồ là
\choice
{$1 : 500$}
{\True $1 : 50\,000$}
{$1 : 5\,000$}
{$1 : 500\,000$}
\loigiai{
Đổi $1{,}5\text{ km} = 150\,000\text{ cm}$.\\
Tỉ lệ bản đồ là $3 : 150\,000 = 1 : 50\,000$.\\
\textbf{Đáp án của câu hỏi là B}
}
\end{ex}

\begin{ex} %[5D7V37]
Một mảnh đất hình vuông có chu vi trên bản đồ tỉ lệ $1 : 500$ là $16\text{ cm}$. Diện tích thực tế của mảnh đất đó là
\choice
{$400\text{ m}^2$}
{$160\text{ m}^2$}
{\True $400\text{ m}^2$}
{$800\text{ m}^2$}
\loigiai{
Cạnh hình vuông trên bản đồ là: $16 : 4 = 4\text{ (cm)}$.\\
Cạnh thực tế của mảnh đất là: $4 \times 500 = 2\,000\text{ cm} = 20\text{ m}$.\\
Diện tích thực tế của mảnh đất là: $20 \times 20 = 400\text{ (m}^2\text{)}$.\\
\textbf{Đáp án của câu hỏi là C}
}
\end{ex}

\begin{ex} %[5D7V37]
Trên bản đồ tỉ lệ $1 : 200$, một mảnh vườn hình chữ nhật có chiều dài $8\text{ cm}$ và chiều rộng $5\text{ cm}$. Diện tích thật của mảnh vườn đó là
\choice
{$160\text{ m}^2$}
{\True $160\text{ m}^2$}
{$1\,600\text{ m}^2$}
{$80\text{ m}^2$}
\loigiai{
Chiều dài thật: $8 \times 200 = 1\,600\text{ cm} = 16\text{ m}$.\\
Chiều rộng thật: $5 \times 200 = 1\,000\text{ cm} = 10\text{ m}$.\\
Diện tích thật: $16 \times 10 = 160\text{ (m}^2\text{)}$.\\
\textbf{Đáp án của câu hỏi là B}
}
\end{ex}

\begin{ex} %[5D7V37]
Trên bản đồ tỉ lệ $1 : 100\,000$, quãng đường từ tỉnh A đến tỉnh B dài $45\text{ cm}$. Một người đi xe máy từ A đến B với vận tốc $30\text{ km/h}$. Thời gian người đó đi hết quãng đường là
\choice
{\True $1{,}5\text{ giờ}$}
{$1\text{ giờ}$}
{$2\text{ giờ}$}
{$2{,}5\text{ giờ}$}
\loigiai{
Quãng đường thực tế: $45 \times 100\,000 = 4\,500\,000\text{ cm} = 45\text{ km}$.\\
Thời gian đi: $45 : 30 = 1{,}5\text{ giờ}$.\\
\textbf{Đáp án của câu hỏi là A}
}
\end{ex}

\begin{ex} %[5D7V37]
Một bể bơi hình chữ nhật có chiều dài thật là $25\text{ m}$, chiều rộng thật là $12\text{ m}$. Khi vẽ trên bản vẽ tỉ lệ $1 : 500$, diện tích hình chữ nhật trên bản vẽ là
\choice
{$6\text{ cm}^2$}
{$24\text{ cm}^2$}
{$60\text{ cm}^2$}
{\True $12\text{ cm}^2$}
\loigiai{
Đổi $25\text{ m} = 2\,500\text{ cm}$; $12\text{ m} = 1\,200\text{ cm}$.\\
Chiều dài trên bản vẽ: $2\,500 : 500 = 5\text{ (cm)}$.\\
Chiều rộng trên bản vẽ: $1\,200 : 500 = 2{,}4\text{ (cm)}$.\\
Diện tích trên bản vẽ: $5 \times 2{,}4 = 12\text{ (cm}^2\text{)}$.\\
\textbf{Đáp án của câu hỏi là D}
}
\end{ex}

\begin{ex} %[5D7V37]
Quãng đường từ Hà Nội đến Hải Phòng dài $105\text{ km}$. Trên một bản đồ, khoảng cách đó là $7\text{ cm}$. Khoảng cách giữa hai điểm C và D trên cùng bản đồ đó là $4\text{ cm}$. Khoảng cách thật giữa C và D là
\choice
{$40\text{ km}$}
{$50\text{ km}$}
{\True $60\text{ km}$}
{$70\text{ km}$}
\loigiai{
Đổi $105\text{ km} = 10\,500\,000\text{ cm}$.\\
Tỉ lệ bản đồ: $7 : 10\,500\,000 = 1 : 1\,500\,000$.\\
Khoảng cách thật giữa C và D là: $4 \times 1\,500\,000 = 6\,000\,000\text{ cm} = 60\text{ km}$.\\
\textbf{Đáp án của câu hỏi là C}
}
\end{ex}

\begin{ex} %[5D7C37]
Một khu đất hình chữ nhật có chu vi thực tế là $240\text{ m}$, chiều rộng bằng $\dfrac{3}{5}$ chiều dài. Khi vẽ khu đất này trên bản đồ với chiều dài $5\text{ cm}$ thì tỉ lệ bản đồ đó là
\choice
{$1 : 500$}
{$1 : 1\,000$}
{$1 : 2\,000$}
{\True $1 : 1\,500$}
\loigiai{
Nửa chu vi thực tế: $240 : 2 = 120\text{ m}$.\\
Chiều dài thực: $120 : (3 + 5) \times 5 = 75\text{ m} = 7\,500\text{ cm}$.\\
Tỉ lệ bản đồ là $5 : 7\,500 = 1 : 1\,500$.\\
\textbf{Đáp án của câu hỏi là D}
}
\end{ex}

\textbf{Phần 2. Trắc nghiệm đúng sai}
\setcounterpageref{ex}{0}

\begin{ex} %[5D7H37]
Trên bản đồ có tỉ lệ $1 : 2\,000\,000$, xét tính đúng sai của các khẳng định sau:
\choiceTF
{\True Đoạn thẳng dài $1\text{ cm}$ trên bản đồ ứng với $20\text{ km}$ ngoài thực tế.}
{Đoạn thẳng dài $5\text{ cm}$ trên bản đồ ứng với $100\text{ m}$ ngoài thực tế.}
{\True Khoảng cách thực tế giữa hai điểm cách nhau $3\text{ cm}$ trên bản đồ là $60\text{ km}$.}
{Nếu khoảng cách thực tế là $80\text{ km}$ thì trên bản đồ vẽ dài $40\text{ cm}$.}
\loigiai{
\begin{itemize}
    \item $1\text{ cm}$ trên bản đồ $= 2\,000\,000\text{ cm} = 20\text{ km}$. Đúng.
    \item $5\text{ cm}$ ứng với $5 \times 20 = 100\text{ km} \neq 100\text{ m}$. Sai.
    \item $3\text{ cm}$ ứng với $3 \times 20 = 60\text{ km}$. Đúng.
    \item $80\text{ km} = 8\,000\,000\text{ cm} : 2\,000\,000 = 4\text{ cm} \neq 40\text{ cm}$. Sai.
\end{itemize}
\begin{itemize}
    \item a) \textbf{Đúng}.
    \item b) \textbf{Sai}.
    \item c) \textbf{Đúng}.
    \item d) \textbf{Sai}.
\end{itemize}
}
\end{ex}

\begin{ex} %[5D7V37]
Một phòng học hình chữ nhật có chiều dài thực tế $8\text{ m}$, chiều rộng $6\text{ m}$. Người ta vẽ mặt bằng phòng học trên bản đồ tỉ lệ $1 : 200$. Xét tính đúng sai của các mệnh đề sau:
\choiceTF
{\True Chiều dài phòng học trên bản đồ là $4\text{ cm}$.}
{\True Chiều rộng phòng học trên bản đồ là $3\text{ cm}$.}
{Chu vi phòng học trên bản đồ là $28\text{ cm}$.}
{\True Diện tích hình chữ nhật biểu diễn phòng học trên bản đồ là $12\text{ cm}^2$.}
\loigiai{
\begin{itemize}
    \item Chiều dài trên bản đồ: $800 : 200 = 4\text{ cm}$.
    \item Chiều rộng trên bản đồ: $600 : 200 = 3\text{ cm}$.
    \item Chu vi trên bản đồ: $(4 + 3) \times 2 = 14\text{ cm} \neq 28\text{ cm}$.
    \item Diện tích trên bản đồ: $4 \times 3 = 12\text{ cm}^2$.
\end{itemize}
\begin{itemize}
    \item a) \textbf{Đúng}.
    \item b) \textbf{Đúng}.
    \item c) \textbf{Sai}.
    \item d) \textbf{Đúng}.
\end{itemize}
}
\end{ex}

\textbf{Phần 3. Trả lời ngắn}
\setcounterpageref{ex}{0}

\begin{ex} %[5D7H37]
Trên bản đồ tỉ lệ $1 : 10\,000$, quãng đường từ cổng trường đến thư viện đo được là $8\text{ cm}$. Tính độ dài thực tế của quãng đường đó theo mét.
\shortans{800}
\loigiai{
\textbf{Đáp số:} $800$\\
$8 \times 10\,000 = 80\,000\text{ cm} = 800\text{ m}$.
}
\end{ex}

\begin{ex} %[5D7H37]
Quãng đường từ nhà An đến công viên dài $2{,}4\text{ km}$. Trên bản đồ tỉ lệ $1 : 40\,000$, khoảng cách này dài bao nhiêu xăng-ti-mét?
\shortans{6}
\loigiai{
\textbf{Đáp số:} $6$\\
$2{,}4\text{ km} = 240\,000\text{ cm} \Rightarrow 240\,000 : 40\,000 = 6\text{ cm}$.
}
\end{ex}

\begin{ex} %[5D7V37]
Trên bản đồ tỉ lệ $1 : 2\,500$, một mảnh đất hình vuông có cạnh dài $4\text{ cm}$. Tính diện tích thực tế của mảnh đất đó theo mét vuông.
\shortans{10000}
\loigiai{
\textbf{Đáp số:} $10000$\\
Cạnh thực tế: $4 \times 2\,500 = 10\,000\text{ cm} = 100\text{ m}$.\\
Diện tích thực: $100 \times 100 = 10\,000\text{ m}^2$.
}
\end{ex}

\begin{ex} %[5D7V37]
Một mảnh đất hình chữ nhật có chu vi $180\text{ m}$, chiều rộng bằng $\dfrac{4}{5}$ chiều dài. Khi vẽ lên bản đồ tỉ lệ $1 : 1\,000$, chu vi hình chữ nhật đó là bao nhiêu xăng-ti-mét?
\shortans{18}
\loigiai{
\textbf{Đáp số:} $18$\\
Đổi $180\text{ m} = 18\,000\text{ cm} \Rightarrow$ Chu vi trên bản đồ: $18\,000 : 1\,000 = 18\text{ cm}$.
}
\end{ex}

\begin{ex} %[5D7V37]
Một công viên hình chữ nhật có chiều dài thực tế $120\text{ m}$, chiều rộng $80\text{ m}$. Trên bản đồ tỉ lệ $1 : 2\,000$, diện tích công viên đó là bao nhiêu xăng-ti-mét vuông?
\shortans{24}
\loigiai{
\textbf{Đáp số:} $24$\\
Chiều dài trên bản đồ: $12\,000 : 2\,000 = 6\text{ cm}$.\\
Chiều rộng trên bản đồ: $8\,000 : 2\,000 = 4\text{ cm}$.\\
Diện tích trên bản đồ: $6 \times 4 = 24\text{ cm}^2$.
}
\end{ex}

\begin{ex} %[5D7C37]
Khoảng cách giữa hai hòn đảo trên bản đồ tỉ lệ $1 : 500\,000$ là $7\text{ cm}$. Một tàu thủy đi từ đảo này sang đảo kia với vận tốc $28\text{ km/h}$. Hỏi tàu thủy đi hết bao nhiêu phút?
\shortans{75}
\loigiai{
\textbf{Đáp số:} $75$\\
Khoảng cách thực tế: $7 \times 500\,000 = 3\,500\,000\text{ cm} = 35\text{ km}$.\\
Thời gian đi: $35 : 28 = 1{,}25\text{ giờ} = 75\text{ phút}$.
}
\end{ex}
